\section{Discussion \& Limitations}
\label{ch:discussion}

% TODO: full prose pass. Below are the load-bearing findings to fold into the discussion narrative.

\subsection{What the Agentic System Can and Cannot Do}

% TODO: synthesise from RQ1/RQ2/RQ3 findings — detection ceiling lifts with scale, replacement quality saturates within the open-weights tier; agentic decomposition's value is in failure traceability and structured ablation rather than in raw automatic metrics.

\subsection{Where Metaphors Fundamentally Resist Literalisation}

% TODO: discuss the H3 evidence: idioms have constrained, conventionalised meanings that map cleanly to literal paraphrases; metaphors --- particularly novel literary metaphors and conceptual metaphors with abstract source--target mappings --- often encode information that no single literal restatement can preserve. Connect to \cite{Neidlein2020} on conventionality stratification.

\subsection{Limitations}

\paragraph{Conceptual limits.} Metaphors encoding meaning that cannot be fully literalised; replacement oversimplifying nuance. The choice to scope the human-evaluation survey to idioms only (Section~\ref{sec:eval-methodology}) reflects this limit pragmatically: the rubric triad of grammaticality, meaning preservation, and simplicity is well-defined for the figurative-to-literal mapping in the SemEval idiom set, but breaks down for conventional metaphors that may not require any rewrite at all.

\paragraph{Evaluation limits.} The subjectivity of human evaluation, mitigated by per-annotator z-score standardisation~\cite{Graham2013} and the assessor-intrinsic QC techniques of bad-references and repeat-pairs~\cite{Graham2017}, remains an inherent property of the task. There is no single ground-truth literal paraphrase for any given idiom, and BLEU's structural misalignment with the task (Section~\ref{sec:rq3}) means the survey's rubric-triad ratings are the load-bearing evidence for replacement quality. The \texttt{n}=17 respondent pool, while sized for $\geq 3$ ratings per item, is at the lower end of what Graham 2013 found stable for per-annotator standardisation; this is documented as a deliberate departure in Section~\ref{sec:eval-methodology}.

\paragraph{Model dependence.} \emph{Findings observed under one model do not necessarily transfer to another, and may invert.} Three instances of this in the thesis, in increasing order of severity:

\textit{(1) v2 chain-of-thought is a no-op on Llama-3.1-8B.} A development-time A/B between prompt v1 (minimal stub) and prompt v2 (chain-of-thought) on \texttt{google/gemini-3-flash-preview} showed $+11$pp BLEU and $+1.6$pp BERTScore F1 in v2's favour on SemEval idiom replacement. On the open-weights deliverable system (Llama-3.1-8B-Instruct) the same A/B produces statistically indistinguishable results: BLEU 0.397 vs.\ 0.400, BERTScore F1 0.914 vs.\ 0.915, sentence F1 0.480 vs.\ 0.492 (Section~\ref{sec:rq3}). The most parsimonious explanation is that the canonical Llama Instruct chat template internalises something close to the v2 scaffold by default, so the explicit prompt-level scaffolding is redundant; the closed-source large model lacks this internal scaffolding and benefits from being told.

\textit{(2) The single-call CoT detection lift inverts at scale, confirmed on both phenomena.} On Llama-3.1-8B-Instruct, the dtr v2 prompt's structured-output schema gives a $+12$pp sentence-F1 lift over the model's own detection-only baseline on VU metaphor (0.430 $\to$ 0.554) and a $+22$pp lift on SemEval idiom (0.255 $\to$ 0.492) --- the auxiliary-replacement field forces the model to commit detected spans alongside a paraphrase, scaffolding detection at small scale. On Llama-3.3-70B-Instruct-AWQ, the same prompt gives a $-28$pp sentence-F1 \textit{loss} from detection-only on VU metaphor (0.635 $\to$ 0.351) and a $-5$pp loss on SemEval idiom (0.435 $\to$ 0.383). Both datasets, same direction, same prompt, no other change in conditions: $+12$ to $+22$pp at 8B, $-5$ to $-28$pp at 70B. Mechanism is consistent across phenomena --- the 70B's token-precision climbs in dtr v2 (more selective flagging) while recall drops (over-commits to ``no figurative content'') under the schema commitment. The 70B is calibrated enough at the detection task that the schema constraint, far from scaffolding the decision, distorts it.

\textit{(3) Pipeline.detect $\equiv$ standalone detection prompt to three decimal places, in all 4 of 4 (model $\times$ dataset) cells.} On the deterministic 200-example slice, detection F1 is identical under standalone detection-only and the 3-step pipeline across both scales and both phenomena: VU metaphor 0.430/0.430 at 8B and 0.635/0.635 at 70B-AWQ; SemEval idiom 0.255/0.255 at 8B and 0.435/0.435 at 70B. All seven token-, span-, and sentence-level detection metrics match across runs to three decimals --- not just sentence F1. The single-call CoT's detection lift over both these conditions is therefore attributable to the structured-output schema's commitment scaffold, not to chain-of-thought reasoning content per se. This rules out the most charitable framing of the v1-vs-v2 finding above (``maybe v2 helps in some other way at 8B'') and pins the scale-inversion finding in (2) to the schema mechanism specifically. It also implies that the pipeline architecture is \textit{scale-invariant for detection}: the detection step in the pipeline gets the same F1 the detection prompt would have gotten alone, regardless of whether the model is 8B or 70B; whatever the pipeline architecture buys is in the replacement step's quality, not in detection sharpness. Replacement-side metrics on SemEval idiom support this read: the +3pp BLEU advantage pipeline shows over CoT-single-call at 8B (0.430 vs.\ 0.400) shrinks to zero at 70B (0.414 vs.\ 0.423) --- the pipeline's BLEU advantage is concentrated at small scale, partially attributable to the ``no-detection $\to$ return-input'' fall-through (input shares $\sim$85\% of tokens with gold), and washes out at larger scale where detection succeeds more often.

These three findings cohere into a single methodological claim: \emph{prompt-engineering wins published on hosted closed-source models, or validated only on small open-weights, give an unreliable read on what helps the deliverable system at production scale.} The risk of taking such findings at face value is publishing a system that performs well in development against the closed-source iterator but does not actually benefit from those design choices in the deliverable --- and, in the most adversarial case, performs \textit{worse} because the prompt scaffolds were calibrated to the smaller model's failure modes. Tian~2024 (TSI), Gao~2025 (PASS), and Badran~2025 (the sequential SemEval Task~1A pipeline) all report architectural wins on hosted closed-source backbones; the thesis's pipeline ablation is a partial empirical re-validation against an open-weights deliverable, and finds the wins partially preserved on detection-replacement decoupling but \textit{not} preserved as detection improvements at scale.

\paragraph{Scale dependence (within family).} The within-family scale comparison (Llama-3.1-8B vs.\ Llama-3.3-70B-Instruct-AWQ) on the 30-source human-evaluation pool yields a clean dissociation: detection sentence F1 lifts from 0.609 to 0.767 (+15.8pp) and span F1 from 0.282 to 0.578 (+29.6pp), but replacement-side BLEU and BERTScore are essentially flat. Scale at this size class lifts the model's ability to \emph{find} the figurative expression but not to \emph{rewrite} it. The 70B comparator is therefore most defensible as a detection-side ceiling for the deliverable rather than a replacement-side competitor.

\paragraph{Scope limits.} Sentence- and MWE-level focus throughout; no discourse-wide figurative-language tracking; primarily English (the FACILE umbrella's Spanish work~\cite{Surez-Figueroa2024} is the natural extension target). The replacement evaluation covers idioms only; metaphor replacement is assessed through detection-side metrics on detect-then-replace runs, not through human-rated rubric scoring.

\subsection{Future Work}

\begin{itemize}
    \item \textbf{Mixture-of-experts evaluation}: Llama 4 (Scout, Maverick) was excluded from this work because the smallest variant exceeds the 40 GB A100's VRAM. With access to A100 80 GB or H100, the within-family scale comparison could be extended into a within-architecture-class comparison, testing whether the agentic decomposition's value transfers to MoE backbones that have different reasoning-vs-throughput tradeoffs.
    \item \textbf{Decomposition ablations}: the no-explanation and no-self-verification ablations (Section~\ref{sec:rq2}) are deferred from this thesis to keep the human-evaluation survey one-shot; isolating the contribution of each agentic step is the natural next study.
    \item \textbf{Curriculum-style data augmentation}: \cite{Jia2024} reaches SOTA on metaphor detection with single-step inference at test time after curriculum-augmented fine-tuning. If the open-weights detection ceiling at 8B (sentence F1 0.430 on VU metaphor detection) is judged insufficient for downstream E2R deployment, CDA-style fine-tuning is the next intervention to try before declaring the open-weights tier insufficient.
    \item \textbf{Spanish FACILE integration}: the per-guideline-module pattern of the UPM E2R cluster~\cite{Surez-Figueroa2024,Surez-Figueroa2023} is language-agnostic at the architectural level. Porting the figurative-language replacement module to Spanish is contingent on a Spanish gold-paraphrase dataset; such a dataset does not currently exist for idioms, so the immediate work is dataset creation rather than model adaptation.
\end{itemize}
